This work sits at the intersection of several adjacent areas. The first is multi-agent planning and orchestration, where the main concern is how to break down work, route subtasks, and coordinate model or tool execution. Representative reasoning-and-acting approaches such as ReAct show the value of interleaving reasoning traces with actions \cite{yao2023react}. That literature is useful for execution strategy, but it often assumes that the main unit of coordination is already given and that success can be measured primarily within the digital action surface. Our concern is earlier and broader: what durable object should carry the work itself, and how should the system behave when the work cannot be reduced to callable digital steps.

The second adjacent area is workflow orchestration and business process management. Those systems contribute important ideas around explicit state, approval gates, and traceability. Workflow-pattern research and artifact-centric process modeling are especially relevant because they emphasize durable state progression and the importance of data-bearing process objects \cite{vanderaalst2003workflow,cohn2009business}. However, many workflow systems are optimized for pre-scoped internal processes rather than ambiguous external demand that must first be clarified, matched, and commercially bounded. The request-rooted model is closer to a commerce-bearing workflow thread than to a predefined internal pipeline.

The third adjacent area is event-sourced and ledger-style application design. Our use of \texttt{RequestEvent} as append-only history inherits the value of immutable audit trails, replay, and projection rebuilds. The novelty here is not event sourcing alone, but its combination with a work-commerce object model that also distinguishes durable proof from ephemeral runtime signals.

The fourth adjacent area is human-in-the-loop and embodied AI. Grounded planning work such as SayCan argues that language plans should be filtered through executable affordances rather than language plausibility alone \cite{ahn2022saycan}. Benchmarks such as ALFRED and TEACh show how difficult embodied instruction following remains even in simulated or constrained environments \cite{shridhar2020alfred,padmakumar2021teach}. Our contribution differs in emphasis. Rather than focusing only on inserting humans into execution, we focus on keeping the requirement for human presence, local access, or physical verification explicit inside the durable request model so that planning cannot erase it.

Finally, the work is adjacent to labor marketplaces and external work platforms. Those systems validate the need for supply discovery, trust, and payment rails, but they often begin from listings or pre-scoped postings. Boreal instead starts from ambiguous demand and treats the request as the durable thread through matching, commitment, fulfillment, proof, and payout. The central distinction of this paper is therefore the use of a request-rooted durable model to preserve execution and proof continuity across digital and embodied work.
